\documentclass[10pt,a4paper]{amsart}
\usepackage{amsmath,amssymb,amsthm,graphicx}
\usepackage{verbatim}
\newenvironment{reference}{\comment}{\endcomment}
%\newenvironment{reference}{}{}
\newenvironment{slogan}{\comment}{\endcomment}
\newenvironment{history}{\comment}{\endcomment}

% For commutative diagrams we use Xy-pic
\usepackage[all]{xy}

% We use 2cell for 2-commutative diagrams.
\xyoption{2cell}
% Theorem environments.
%
\theoremstyle{plain}
\newtheorem{theorem}[subsection]{Theorem}
\newtheorem{proposition}[subsection]{Proposition}
\newtheorem{lemma}[subsection]{Lemma}

\theoremstyle{definition}
\newtheorem{definition}[subsection]{Definition}
\newtheorem{example}[subsection]{Example}
\newtheorem{exercise}[subsection]{Exercise}
\newtheorem{situation}[subsection]{Situation}

\theoremstyle{remark}
\newtheorem{remark}[subsection]{Remark}
\newtheorem{remarks}[subsection]{Remarks}

\numberwithin{equation}{subsection}

% Macros
%
\def\lim{\mathop{\mathrm{lim}}\nolimits}
\def\colim{\mathop{\mathrm{colim}}\nolimits}
\def\Spec{\mathop{\mathrm{Spec}}}
\def\Hom{\mathop{\mathrm{Hom}}\nolimits}
\def\Ext{\mathop{\mathrm{Ext}}\nolimits}
\def\SheafHom{\mathop{\mathcal{H}\!\mathit{om}}\nolimits}
\def\SheafExt{\mathop{\mathcal{E}\!\mathit{xt}}\nolimits}
\def\Sch{\mathit{Sch}}
\def\Mor{\mathop{\mathrm{Mor}}\nolimits}
\def\Ob{\mathop{\mathrm{Ob}}\nolimits}
\def\Sh{\mathop{\mathit{Sh}}\nolimits}
\def\NL{\mathop{N\!L}\nolimits}
\def\CH{\mathop{\mathrm{CH}}\nolimits}
\def\proetale{{pro\text{-}\acute{e}tale}}
\def\etale{{\acute{e}tale}}
\def\QCoh{\mathit{QCoh}}
\def\Ker{\mathop{\mathrm{Ker}}}
\def\Im{\mathop{\mathrm{Im}}}
\def\Coker{\mathop{\mathrm{Coker}}}
\def\Coim{\mathop{\mathrm{Coim}}}

% Boxtimes
%
\DeclareMathSymbol{\boxtimes}{\mathbin}{AMSa}{"02}

%
% Macros for moduli stacks/spaces
%
\def\QCohstack{\mathcal{QC}\!\mathit{oh}}
\def\Cohstack{\mathcal{C}\!\mathit{oh}}
\def\Spacesstack{\mathcal{S}\!\mathit{paces}}
\def\Quotfunctor{\mathrm{Quot}}
\def\Hilbfunctor{\mathrm{Hilb}}
\def\Curvesstack{\mathcal{C}\!\mathit{urves}}
\def\Polarizedstack{\mathcal{P}\!\mathit{olarized}}
\def\Complexesstack{\mathcal{C}\!\mathit{omplexes}}
% \Pic is the operator that assigns to X its picard group, usage \Pic(X)
% \Picardstack_{X/B} denotes the Picard stack of X over B
% \Picardfunctor_{X/B} denotes the Picard functor of X over B
\def\Pic{\mathop{\mathrm{Pic}}\nolimits}
\def\Picardstack{\mathcal{P}\!\mathit{ic}}
\def\Picardfunctor{\mathrm{Pic}}
\def\Deformationcategory{\mathcal{D}\!\mathit{ef}}
\usepackage{cleveref}
\crefname{originallemma}{original lemma}{original lemmas}
\Crefname{originallemma}{Original lemma}{Original lemmas}
\makeatletter
\newlabel{lemma-extension-by-zero}{{0G4L}{source}{}{sourcepointer-1}{}}
\newlabel{lemma-extension-by-zero@cref}{{[originallemma][0][]0G4L}{[1][0][]source}{}{}{}}
\makeatother
\crefname{originallemma}{original lemma}{original lemmas}
\Crefname{originallemma}{Original lemma}{Original lemmas}
\makeatletter
\newlabel{lemma-lift-map}{{0G2H}{source}{}{sourcepointer-2}{}}
\newlabel{lemma-lift-map@cref}{{[originallemma][0][]0G2H}{[1][0][]source}{}{}{}}
\makeatother
\crefname{originalproposition}{original proposition}{original propositions}
\Crefname{originalproposition}{Original proposition}{Original propositions}
\makeatletter
\newlabel{perfect-proposition-DCoh}{{0FDB}{source}{}{sourcepointer-3}{}}
\newlabel{perfect-proposition-DCoh@cref}{{[originalproposition][0][]0FDB}{[1][0][]source}{}{}{}}
\makeatother
\crefname{originalsection}{original section}{original sections}
\Crefname{originalsection}{Original section}{Original sections}
\makeatletter
\newlabel{derived-section-cones}{{014D}{source}{}{sourcepointer-4}{}}
\newlabel{derived-section-cones@cref}{{[originalsection][0][]014D}{[1][0][]source}{}{}{}}
\makeatother
\crefname{originalsection}{original section}{original sections}
\Crefname{originalsection}{Original section}{Original sections}
\makeatletter
\newlabel{derived-section-canonical-delta-functor}{{014Z}{source}{}{sourcepointer-5}{}}
\newlabel{derived-section-canonical-delta-functor@cref}{{[originalsection][0][]014Z}{[1][0][]source}{}{}{}}
\makeatother
\crefname{originallemma}{original lemma}{original lemmas}
\Crefname{originallemma}{Original lemma}{Original lemmas}
\makeatletter
\newlabel{derived-lemma-third-isomorphism-triangle}{{014A}{source}{}{sourcepointer-6}{}}
\newlabel{derived-lemma-third-isomorphism-triangle@cref}{{[originallemma][0][]014A}{[1][0][]source}{}{}{}}
\makeatother
\title{A distinguished triangle on an open lifts to a compatible system of distinguished triangles on the ambient scheme}
\author{The Stacks Project Authors}
\date{stacks-project:duality.tex:local-snapshot-20261012; GNU FDL 1.2 or later}
\begin{document}
\maketitle
\noindent\textit{Editorial note.} All delivered mathematical text is the author's own, copied from the pinned Stacks Project chapter file duality.tex; no statement or proof step was rewritten, reordered or removed.\par\medskip
\noindent\textit{Editorial note.} Preamble provenance: the pinned chapter file opens with an input of the shared preamble and carries no preamble text of its own, so the project's single Stacks-wide preamble staged at sources/stacks-duality/preamble.tex is excerpted verbatim under the recipe fixed by LANE\_BRIEF 22.1, as three spans, lines 8-13 (verbatim package and the reference/slogan/history environments), lines 14-18 (Xy-pic setup) and lines 154-217 (the native theorem declarations and mathematical macros). Nothing was added, removed or redefined; the wrapper class and the amsmath/amssymb/amsthm/graphicx base packages are supplied by the assembler.\par\medskip
\noindent\textit{Editorial note.} The counted statement and proof contain no cross reference, no citation and no included file, so the delivered body is self contained.\par\medskip
\noindent\textit{Editorial note.} Canonical location of the counted proposition: Stacks tag 0G4M (sources/stacks-duality/pinned-original-source.tex line 8145).\par\medskip
\noindent\textit{Allowed standard prerequisites (editorial).}\begin{itemize}
\item Existence of Deligne systems and of lifts of bounded coherent objects from an open. Same chapter, Lemma labelled lemma-extension-by-zero part (2), cited at proof lines 8161-8163. Provides (K\_n) and L'.
\item Surjectivity of Hom onto the open for a Deligne system. Same chapter, Lemma labelled lemma-lift-map, cited at proof lines 8163-8166. Lifts the arrow K -> L to a morphism over X.
\item The bounded coherent derived category as the derived category of the abelian category of coherent modules. Derived Categories of Schemes, Proposition labelled perfect-proposition-DCoh, cited at proof lines 8176-8178. Represents the lifted morphism by a morphism of complexes.
\item Cones of complexes and axiom TR3. Derived Categories, Section labelled derived-section-cones and Lemma labelled derived-lemma-third-isomorphism-triangle, cited at proof lines 8188-8200. Show that the cones form the required compatible triangles.
\end{itemize}\medskip
\ifdefined\hypertarget\hypertarget{sourcepointer-1}{}\fi
\noindent\textit{Source reference.} Original lemma 0G4L: sources/stacks-duality/pinned-original-source.tex, source lines 8103--8103. Same-chapter pointer into the author environment lemma at source line 8103; invoked at proof line 8162 of the counted proof. It is retained only as a navigation pointer; every load-bearing step is inside the delivered span.\par\medskip
\ifdefined\hypertarget\hypertarget{sourcepointer-2}{}\fi
\noindent\textit{Source reference.} Original lemma 0G2H: sources/stacks-duality/pinned-original-source.tex, source lines 7920--7920. Same-chapter pointer into the author environment lemma at source line 7920; invoked at proof line 8165 of the counted proof. It is retained only as a navigation pointer; every load-bearing step is inside the delivered span.\par\medskip
\ifdefined\hypertarget\hypertarget{sourcepointer-3}{}\fi
\noindent\textit{Source reference.} Original proposition 0FDB: sources/stacks-chow/tags.txt, source lines 19211--19211. Cross-chapter standard prerequisite located in the project tag list staged in the local snapshot; invoked at proof line 8172 of the counted proof. It is retained only as a navigation pointer; every load-bearing step is inside the delivered span.\par\medskip
\ifdefined\hypertarget\hypertarget{sourcepointer-4}{}\fi
\noindent\textit{Source reference.} Original section 014D: sources/stacks-chow/tags.txt, source lines 1468--1468. Cross-chapter standard prerequisite located in the project tag list staged in the local snapshot; invoked at proof line 8185 of the counted proof. It is retained only as a navigation pointer; every load-bearing step is inside the delivered span.\par\medskip
\ifdefined\hypertarget\hypertarget{sourcepointer-5}{}\fi
\noindent\textit{Source reference.} Original section 014Z: sources/stacks-chow/tags.txt, source lines 1490--1490. Cross-chapter standard prerequisite located in the project tag list staged in the local snapshot; invoked at proof line 8193 of the counted proof. It is retained only as a navigation pointer; every load-bearing step is inside the delivered span.\par\medskip
\ifdefined\hypertarget\hypertarget{sourcepointer-6}{}\fi
\noindent\textit{Source reference.} Original lemma 014A: sources/stacks-chow/tags.txt, source lines 1465--1465. Cross-chapter standard prerequisite located in the project tag list staged in the local snapshot; invoked at proof line 8199 of the counted proof. It is retained only as a navigation pointer; every load-bearing step is inside the delivered span.\par\medskip
\begin{lemma}
\label{lemma-extension-by-zero-triangle}
Let $j : U \to X$ be an open immersion of Noetherian schemes. Let
$$
K \to L \to M \to K[1]
$$
be a distinguished triangle of $D^b_{\textit{Coh}}(\mathcal{O}_U)$.
Then there exists an inverse system of distinguished triangles
$$
K_n \to L_n \to M_n \to K_n[1]
$$
in $D^b_{\textit{Coh}}(\mathcal{O}_X)$ such that $(K_n)$, $(L_n)$, $(M_n)$
are Deligne systems and such that the restriction of these
distinguished triangles to $U$ is isomorphic to the distinguished triangle
we started out with.
\end{lemma}

\begin{proof}
Let $(K_n)$ be as in Lemma \ref{lemma-extension-by-zero} part (2).
Choose an object $L'$ of $D^b_{\textit{Coh}}(\mathcal{O}_X)$
whose restriction to $U$ is $L$ (we can do this as the lemma shows).
By Lemma \ref{lemma-lift-map} we can find an $n$ and a morphism
$K_n \to L'$ on $X$ whose restriction to $U$ is the given arrow
$K \to L$. We conclude there is a morphism $K' \to L'$ of
$D^b_{\textit{Coh}}(\mathcal{O}_X)$ whose restriction to $U$
is the given arrow $K \to L$.

\medskip\noindent
By Derived Categories of Schemes, Proposition \ref{perfect-proposition-DCoh}
we can find a morphism $\alpha^\bullet : \mathcal{F}^\bullet \to
\mathcal{G}^\bullet$ of bounded complexes
of coherent $\mathcal{O}_X$-modules representing $K' \to L'$.
Choose a quasi-coherent sheaf of ideals $\mathcal{I} \subset \mathcal{O}_X$
whose vanishing locus is $X \setminus U$.
Then we let
$K_n = \mathcal{I}^n\mathcal{F}^\bullet$
and
$L_n = \mathcal{I}^n\mathcal{G}^\bullet$.
Observe that $\alpha^\bullet$ induces a morphism of
complexes $\alpha_n^\bullet : \mathcal{I}^n\mathcal{F}^\bullet \to
\mathcal{I}^n\mathcal{G}^\bullet$. From the construction of cones in
Derived Categories, Section \ref{derived-section-cones}
it is clear that
$$
C(\alpha_n)^\bullet = \mathcal{I}^nC(\alpha^\bullet)
$$
and hence we can set $M_n = C(\alpha_n)^\bullet$. Namely, we
have a compatible system of distinguished triangles
(see discussion in Derived Categories, Section
\ref{derived-section-canonical-delta-functor})
$$
K_n \to L_n \to M_n \to K_n[1]
$$
whose restriction to $U$ is isomorphic to the distinguished 
triangle we started out with by axiom TR3 and Derived Categories,
Lemma \ref{derived-lemma-third-isomorphism-triangle}.
\end{proof}

\end{document}
